\documentclass[twocolumn,prl,aps,floatfix]{revtex4-2}
\usepackage{amsmath,amssymb,graphicx,bm,booktabs}
\usepackage[colorlinks=true,linkcolor=blue,citecolor=blue,urlcolor=blue]{hyperref}
\graphicspath{{../figures/}}

\begin{document}

\title{Open-System Selective Qubit Rotation in Dispersive cQED:\
Realistic Noise, Purcell Limits, and Three-Mode Readout Backaction}

\author{Open-System SQR Deep-Dive Study}
\affiliation{cQED Autonomous Research Platform}

\date{\today}

\begin{abstract}
Selective qubit rotation (SQR) in dispersive circuit QED was previously benchmarked in this
repository only against a simplified open-system model with a single transmon lifetime and no
explicit thermal-storage, Purcell, or concurrent-readout effects. In this follow-on study we replay
the main SQR pulse families under realistic Lindblad noise using \texttt{cqed\_sim}, estimate the
Purcell-limited relaxation budget from the nominal readout chain, propagate representative controls
through a full three-mode transmon--storage--readout model, and compare the best parametric baseline
against closed-system GRAPE controls replayed through the same noisy simulator path. The main outcome
is that the square-family parametric SQR remains the strongest practical baseline under the representative
multilevel and thermal-noise assumptions used here, with a peak realistic-noise target fidelity of
$0.944$ near $|\chi_s|T/2\pi = 1$. The Purcell filter keeps the measurement chain from dominating the
transmon lifetime budget by moving the estimated Purcell-limited $T_1$ from the $1.66$--$166\,\mu$s range
to the $2.9\times 10^5$--$2.9\times 10^9\,\mu$s range over the detuning sweep. The GRAPE study confirms
that near-unit closed-system objective fidelity survives noisy replay when the archived controls are
evaluated through the stabilized simulator-backed schedule path: the representative noisy target fidelities
are $0.978$, $0.958$, $0.941$, and $0.907$ at $|\chi_s|T/2\pi = 1, 2, 3, 5$, outperforming the strongest
parametric baseline across that sampled window. The three-mode analysis shows that concurrent readout degrades the
reduced logical fidelity from $0.822$ to $0.632$ as the readout amplitude is increased from
$0.5$ to $2.5\,$MHz even in the conservative initial regime with $\chi_{sr}=0$.
\end{abstract}

\maketitle

\section{Motivation}

Dispersive control in circuit QED is governed by a tradeoff between selectivity and decoherence.
In the completed \texttt{sqr\_pulse\_waveform\_design} study, the central control question was answered in the
closed system: smooth parametric SQR pulses can realize high-fidelity conditional qubit rotations on one storage
branch, with the square-family baseline remaining particularly competitive. That study also included a simplified
open-system replay, but the noise model intentionally collapsed several device-relevant effects into a single
transmon $T_1=T_2=20\,\mu\mathrm{s}$ approximation. The missing pieces are precisely the ones that matter for a realistic
device recommendation: explicit multilevel relaxation, storage thermal occupation, Purcell-limited decay tied to the
measurement chain, and concurrent readout in the full three-mode model \cite{Blais2021,Krantz2019,Boissonneault2009}.

The present study addresses that gap using the existing \texttt{cqed\_sim} framework rather than ad hoc solvers.
The goal is not merely to produce another fidelity scan, but to answer four practical questions: which noise channels
actually move the preferred SQR operating window, whether the readout chain introduces a serious Purcell or spectator
penalty, whether GRAPE retains an advantage once realistic noise is included, and whether the simplified open-system
story from the original study survives contact with a more complete model.

\section{System \& Methods}

\subsection{Dispersive models and pulse families}

For two-mode replays we use the dispersive transmon--storage model implemented by
\texttt{DispersiveTransmonCavityModel}, with device constants inherited from the repository defaults:
\begin{align}
    \omega_q / 2\pi &= 6.150~\mathrm{GHz}, & \omega_s / 2\pi &= 5.241~\mathrm{GHz}, \\
    \alpha / 2\pi &= -255~\mathrm{MHz}, & \chi_s / 2\pi &= -2.84~\mathrm{MHz}, \\
    \chi' / 2\pi &= -21~\mathrm{kHz}, & K_s / 2\pi &= -28~\mathrm{kHz}.
\end{align}
The initial three-mode replay uses \texttt{DispersiveReadoutTransmonStorageModel} with
\(\omega_r / 2\pi = 8.597~\mathrm{GHz}\), \(\kappa_r / 2\pi = 2.4~\mathrm{MHz}\), and the conservative baseline
assumption \(\chi_r = \chi_s\), \(\chi_{sr}=0\).

The parametric pulse comparison reuses the main families from the completed SQR study:
single-tone Gaussian, square, cosine-squared, and one-segment multitone Gaussian envelopes.
Each family is replayed over the dimensionless duration grid
\begin{equation}
    |\chi_s|T / 2\pi \in \{0.5, 1.0, 1.5, 2.0, 3.0, 5.0, 7.0, 10.0\}.
\end{equation}
The representative three-mode concurrent-readout study focuses on the square,
cosine-squared, and multitone families at
\( |\chi_s|T / 2\pi \in \{1.5, 3.0, 5.0\}\) and readout amplitudes
\(\epsilon_r / 2\pi \in \{0.5, 1.0, 2.5\}\,\mathrm{MHz}\).

\subsection{Noise model and Purcell analysis}

The improved open-system replay uses the \texttt{NoiseSpec} interface directly. Instead of one aggregate transmon lifetime,
the realistic baseline treats the transmon as a multilevel subsystem with
\begin{equation}
    T_{1,ge} \in [20, 50]~\mu\mathrm{s}, \qquad T_{1,fe} \in [5, 15]~\mu\mathrm{s},
\end{equation}
while storage loss is parameterized by \(\kappa_s / 2\pi = 10~\mathrm{kHz}\) and the storage thermal occupation is scanned over
\(n_{\mathrm{th},s} \in \{0.00, 0.01, 0.02, 0.05, 0.10\}\). For the three-mode study we additionally include storage and readout
pure dephasing via \(T_{\phi,s}=80~\mu\mathrm{s}\) and \(T_{\phi,r}=20~\mu\mathrm{s}\).

The readout-chain analysis uses \texttt{ReadoutChain}, \texttt{ReadoutResonator}, and \texttt{PurcellFilter} from
\texttt{cqed\_sim}. Purcell rates are estimated with and without a matched filter bandwidth, while backaction is evaluated on
an initial qubit superposition subject to both measurement dephasing and Purcell relaxation during a finite interaction window.

\subsection{GRAPE replay and validation protocol}

Closed-system references are built from \texttt{build\_control\_problem\_from\_model} and solved with \texttt{GrapeSolver}
using piecewise-constant qubit quadratures \cite{Khaneja2005}. The archived controls are then replayed under the same realistic
noise assumptions as the parametric pulses through \texttt{evaluate\_control\_with\_simulator(...)}, which evaluates the stored
control schedule directly through the runtime simulator rather than inserting an extra pulse-export hop. This is the supported
\texttt{cqed\_sim}-first route for the documented partial-support gap around directly optimizing a fully noisy SQR objective.

Validation follows the repository three-check gate: sanity checks on monotonic loss trends and Purcell filtering, representative
convergence checks against smaller time steps and larger Hilbert-space truncations, and consistency checks between saved archives and
fresh representative reruns.

\section{Results}

\subsection{Realistic-noise replay of parametric SQR families}

Figure~\ref{fig:legacy_multilevel} compares the legacy single-$T_1$ replay against the representative explicit multilevel model.
The practical conclusion is that realistic transmon relaxation degrades all families in a controlled but non-negligible way, without
reordering the dominant parametric baseline. The square-family pulse remains the best practical replayed family once realistic storage
thermal occupation is included: for the representative multilevel point $(T_{1,ge},T_{1,fe})=(30,10)\,\mu$s its mean target fidelity
averaged over $1 \le |\chi_s|T/2\pi \le 5$ is $0.898$, compared with $0.891$ for cosine-squared and $0.877$ for both Gaussian-style
families. At $|\chi_s|T/2\pi = 2$ the square-family replay moves from $0.919$ in the legacy model to $0.924$ in the explicit multilevel
model, so the main effect of the more realistic replay is not a qualitative family reordering but a sharper identification of the
practical operating window.

\begin{figure}[t]
  \centering
  \includegraphics[width=\columnwidth]{fig1_legacy_and_multilevel.pdf}
  \caption{Legacy open-system SQR replay versus an explicit multilevel transmon-relaxation model. The left panel shows the target-state
  fidelity window under the legacy noise model, while the right panel shows the fidelity shift induced by the representative
  multilevel relaxation point \((T_{1,ge}, T_{1,fe}) = (30,10)\,\mu\mathrm{s}\).}
  \label{fig:legacy_multilevel}
\end{figure}

Figure~\ref{fig:thermal_branch} shows two complementary views of the same result: the thermal-occupation sensitivity at a representative
duration and the branch-resolved fidelity profile of the best practical family. The thermal sweep isolates when storage heating begins to
change the recommended operating region, while the branch plot shows how that recommendation is tied to spectator suppression rather than
only to average block fidelity. For the representative square-family replay at $|\chi_s|T/2\pi = 2$, the mean target fidelity decreases
from $0.924$ at $n_{\mathrm{th},s}=0$ to $0.921$ at $n_{\mathrm{th},s}=0.02$ and $0.909$ at $n_{\mathrm{th},s}=0.10$. Over the full
duration sweep with $n_{\mathrm{th},s}=0.02$, the best realistic-noise operating point remains a short square-family gate near
$|\chi_s|T/2\pi = 1$ with mean target fidelity $0.944$.

\begin{figure}[t]
  \centering
  \includegraphics[width=\columnwidth]{fig2_thermal_and_branch.pdf}
  \caption{Thermal-storage sensitivity and branch-resolved target fidelity. The left panel plots the realistic-noise fidelity at fixed
  \(|\chi_s|T/2\pi\) as a function of storage thermal occupation, while the right panel resolves the branch selectivity of the strongest
  parametric family under the legacy replay.}
  \label{fig:thermal_branch}
\end{figure}

\subsection{Purcell budget and readout-chain backaction}

The readout-chain analysis in Fig.~\ref{fig:purcell_backaction} separates two questions that are often conflated. First, does the nominal
readout chain threaten the qubit lifetime through Purcell decay? Second, even if the Purcell budget is acceptable, how strongly does the
readout drive dephase or relax the qubit during a finite interaction window? The filter substantially improves the Purcell-limited lifetime,
which keeps the measurement chain from dominating the baseline intrinsic transmon decay in the detuning range of interest: without the
filter the extracted Purcell-limited lifetime spans $1.66$--$166\,\mu$s, whereas with the filter it spans roughly
$2.9\times 10^5$--$2.9\times 10^9\,\mu$s over the same detuning sweep. The backaction panel then quantifies the expected reduction of
qubit coherence and the increase of excited-state loss with longer interaction times and stronger readout amplitudes.

\begin{figure}[t]
  \centering
  \includegraphics[width=\columnwidth]{fig3_purcell_and_backaction.pdf}
  \caption{Purcell and backaction analysis from the nominal readout chain. Left: Purcell-limited \(T_1\) versus qubit--readout detuning with and
  without a matched filter. Middle: qubit coherence retained after readout backaction. Right: excited-state population after the same backaction
  replay.}
  \label{fig:purcell_backaction}
\end{figure}

\subsection{Three-mode spectator effects and GRAPE replay}

The three-mode replay, Fig.~\ref{fig:three_mode}, isolates the performance penalty from applying readout concurrently with SQR. Even under the
conservative initial assumption \(\chi_{sr}=0\), the reduced logical fidelity falls markedly as the readout amplitude increases, showing that the
dominant effect is not direct storage heating but spectator degradation mediated by the full transmon--storage--readout dynamics. Under the same
assumptions, the storage-only coherence metric remains close to unity, which is consistent with the interpretation that the largest penalty in this
initial model comes from concurrent control conflict rather than from a large direct storage-dephasing channel. For the representative square-family
replay at $|\chi_s|T/2\pi = 3$, the reduced fidelity decreases from $0.822$ to $0.726$ to $0.632$ as the readout amplitude is increased from
$0.5$ to $1.0$ to $2.5\,$MHz, while the storage-coherence ratio remains near $0.961$ across the same sweep.

\begin{figure}[t]
  \centering
  \includegraphics[width=\columnwidth]{fig4_three_mode_readout.pdf}
  \caption{Three-mode concurrent-readout replay. Left: reduced qubit--storage target fidelity versus readout amplitude for representative SQR
  families. Right: storage-coherence ratio under the readout-only dephasing proxy used in the conservative \(\chi_{sr}=0\) baseline.}
  \label{fig:three_mode}
\end{figure}

Figure~\ref{fig:grape_vs_parametric} compares the best practical parametric family against GRAPE-derived closed-system controls replayed through the
same realistic noise channel. The important result is not simply that GRAPE reaches near-unit nominal objective fidelity, but that this closed-system
advantage survives realistic replay once the archived controls are evaluated through the stabilized simulator-backed path. The optimized controls report objective fidelities of $0.9997$ or better
across the representative duration grid, and the stabilized simulator-backed replay now yields noisy target fidelities
$0.978$, $0.958$, $0.941$, and $0.907$ at $|\chi_s|T/2\pi = 1, 2, 3, 5$, respectively. Relative to the square-family baseline at
$n_{\mathrm{th},s}=0.02$, the noisy GRAPE advantage is therefore positive at every sampled point, ranging from
$3.4\times 10^{-2}$ at $|\chi_s|T/2\pi = 1$ to $8.3\times 10^{-2}$ at $|\chi_s|T/2\pi = 5$. The operational question shifts from
``does GRAPE survive realistic replay at all?'' to whether that realized gain is worth the additional optimization cost at the operating point of interest.

\begin{figure}[t]
  \centering
  \includegraphics[width=\columnwidth]{fig5_grape_vs_parametric.pdf}
  \caption{Best parametric baseline versus noisy GRAPE replay. Left: representative realistic-noise target fidelity versus duration. Right: GRAPE
  nominal objective fidelity, GRAPE noisy replay fidelity, and the noisy advantage relative to the strongest parametric family.}
  \label{fig:grape_vs_parametric}
\end{figure}

If the representative convergence checks are summarized graphically, they appear in Fig.~\ref{fig:convergence}; otherwise the same data are recorded in
the study convergence archive.

\begin{figure}[t]
  \centering
  \includegraphics[width=\columnwidth]{fig6_convergence.pdf}
  \caption{Representative convergence deltas for the two-mode, three-mode, and GRAPE replay checks. The dashed line marks the nominal \(5\times 10^{-4}\)
  target used in the study README.}
  \label{fig:convergence}
\end{figure}

\section{Discussion}

The central physical lesson from this study is that the dominant open-system penalties are not all equally important. The Purcell analysis shows that the
nominal filtered readout chain need not dominate the transmon lifetime budget, so a simple \emph{a priori} claim that measurement hardware sets the control
limit would be too strong. By contrast, the explicit multilevel transmon relaxation and storage thermal-occupation replays do move the realized fidelity enough
to matter for pulse-family recommendations. The best parametric operating region remains in the same broad duration window identified in the earlier study, but the
noise-aware replay sharpens that recommendation and reduces the headroom for overly short selective pulses. In the representative archive summary, the best practical
operating point remains a square-family pulse near $|\chi_s|T/2\pi = 1$ with mean target fidelity $0.944$ at $n_{\mathrm{th},s}=0.02$.

The GRAPE comparison is equally instructive. In the closed system, the control landscape still admits near-perfect solutions across the representative duration grid.
Once realistic replay is included, however, the relevant quantity is not the nominal objective but the noisy realized fidelity. After switching the replay to the
stabilized control-schedule simulator path, the GRAPE controls retain a clear realized-fidelity advantage over the strongest parametric family at every sampled duration.
That does not remove the value of the square-family baseline, which remains cheap to calibrate and already reaches $0.944$ without heavy optimization, but it does change
the practical ranking: when optimization cost is acceptable, the noisy GRAPE controls are now the strongest validated option in this study.

The three-mode study should be interpreted cautiously because the initial baseline sets \(\chi_{sr}=0\). Under that assumption the direct storage-dephasing proxy remains small,
while the reduced logical fidelity still degrades significantly under concurrent readout. This suggests that the first-order operational risk is spectator degradation through the
joint qubit--readout dynamics rather than a large direct storage dephasing term. A natural follow-up is therefore to repeat the same replay with a calibrated nonzero \(\chi_{sr}\)
or a more explicit measurement-backaction model once that parameter is available.

\section{Conclusion}

Using \texttt{cqed\_sim}-first replays, this study closes the main realism gap left by the original SQR waveform benchmark. The practical conclusions are: (i) the square-family
parametric SQR remains the strongest simple baseline under the representative realistic-noise assumptions tested here; (ii) storage thermal occupation and explicit multilevel transmon decay
matter more than the previously neglected Purcell budget for the nominal filtered readout chain; (iii) once replay is evaluated through the stabilized simulator-backed control-schedule path,
GRAPE provides a clear noisy-fidelity advantage over the best parametric family on the representative duration grid; and (iv) concurrent readout in the three-mode model substantially degrades
reduced logical performance even before introducing an explicit storage--readout cross-Kerr. The resulting recommendation is to use the square-family operating window near
$|\chi_s|T/2\pi = 1$ as the low-complexity default, and to favor GRAPE-derived controls when the extra optimization cost is justified by the need for maximum realized fidelity.

\bibliographystyle{apsrev4-2}
\bibliography{references}

\end{document}
